\documentclass[10pt,a4paper]{amsart}
\usepackage[margin=19mm]{geometry}
\usepackage{amsmath,amssymb,mathtools}
\usepackage[scr=rsfs,cal=euler]{mathalfa}
\usepackage{xfrac}
\usepackage[unicode,hidelinks,bookmarks=false]{hyperref}
\newcommand{\ie}{i.e.\ }
\newcommand{\cf}{cf.\ }
\newcommand{\resp}{respectively\ }
\newcommand{\Subsection}{\subsection}
\providecommand{\nobreakdash}{\nobreak\hbox{-}}
\setlength{\emergencystretch}{2em}
\multlinegap0pt
\usepackage{bm}
\usepackage{bbm}
\usepackage{dsfont}
\usepackage{xspace}
\usepackage{cancel}
\usepackage{mathtools}
\usepackage{amsthm}
\makeatletter
\@ifundefined{theorem}{\newtheorem{theorem}{Theorem}}{}
\@ifundefined{definition}{\newtheorem{definition}[theorem]{Definition}}{}
\@ifundefined{proposition}{\newtheorem{proposition}[theorem]{Proposition}}{}
\makeatother
\title[Twisted Gelfand-Ponomarev modules]{Twisted Gelfand-Ponomarev modules}
\author{Joseph Muller, Chia-Fu Yu}
\date{Source: arXiv:2603.12116v2 (CC BY 4.0)}
\begin{document}
\maketitle

\noindent\emph{Source and licence.} Twisted Gelfand-Ponomarev modules. Version arXiv:2603.12116v2 (CC BY 4.0), distributed under the Creative Commons Attribution 4.0 International licence, as recorded in the arXiv licence metadata for this exact version and shown on the abstract page https://arxiv.org/abs/2603.12116v2. Full rights notice: licenses/arxiv-2603.12116-CC-BY-4.0.txt.\par\smallskip

\noindent\textit{Editorial scope.} Counted unit: the Proposition whose source label is \texttt{Gr1stIsOfTheFirstKind} in the article (source0.tex lines 1663--1665, source label \texttt{Gr1stIsOfTheFirstKind}), delivered together with the article's own complete original proof of it (source0.tex lines 1667--1735, one \texttt{proof} environment of 69 lines). No mathematics is written, repaired or re-derived: statement and proof are the article's own text line for line. Required context retained uncounted: DefinitionModuleAttachedToKraftQuiver (lines 1218--1231). Every definition, display and label of those lines is retained unchanged. Omitted from this package and not redistributed: 1--1217, 1232--1662, 1666--1666, 1736--2638. Nothing inside a retained excerpt is omitted or altered: every retained body is delivered exactly as the article's lines are written, so no omission receipt of any kind accompanies this package. Citations: the retained excerpts carry no citation command at all, so no bibliography entry of the article is transcribed and no reference is redistributed. Reference adaptation: none was needed and none was made; every reference inside the delivered unit resolves inside it, so no unresolved reference or citation is produced. Printed slips kept literal: no printed word, symbol, hypothesis or number of the counted statement or of its proof was altered. Layout and class: the article's own class is \texttt{amsart} with packages including etex, pstricks, epsfig, graphicx, color, geometry, xy, amssymb, amscd, cite, fullpage, marvosym, url, comment; the wrapper is the lane's pinned \texttt{amsart} editorial wrapper at 10pt on A4, which restates the theorem-like environments the retained text needs and adds the article's own macro definitions that the retained text needs (none), with extra wrapper packages (bm, bbm, dsfont, xspace, cancel, mathtools). The delivered numbering is the unit's own. Assets: no retained span contains a figure environment, a graphics-inclusion command, a PDF-page inclusion, a table or a TikZ picture; no image file is copied into this package, the package declares no asset dependency and no image file is redistributed with it. Licence: version arXiv:2603.12116v2 of this article is distributed under the Creative Commons Attribution 4.0 International licence, as recorded in the arXiv licence metadata for this exact version and shown on the abstract page https://arxiv.org/abs/2603.12116v2. A full rights notice is delivered at licenses/arxiv-2603.12116-CC-BY-4.0.txt. Characters: every retained excerpt is pure ASCII, so the delivered engine, XeLaTeX with its default Latin Modern text font, reports no missing character.
\begin{definition}\label{DefinitionModuleAttachedToKraftQuiver}
    Let $(U,\rho)$ be a $(\sigma,\tau)$-linear representation of a finite directed graph $\Gamma$ labeled by $\mathcal A$. The \textit{module attached to $(U,\rho)$} is the left $K\langle F,V\rangle_{\sigma,\tau}$-module defined by 
    \begin{equation*}
        M(\Gamma,U,\rho) := \bigoplus_{v \in \mathcal V} U_v,
    \end{equation*}
    together with the operators $F$ and $V$ whose restrictions to a given summand $U_v$ are
    \begin{equation*}
        F_{|U_v} = \bigoplus_{\substack{E = (v,v') \in \mathcal E\\ \mathrm{label}(E)  = F}} \rho_E,
    \end{equation*}
    and 
    \begin{equation*}
        V_{|U_v} = \bigoplus_{\substack{E = (v,v') \in \mathcal E\\ \mathrm{label}(E)  = V}} \rho_E.
    \end{equation*}
\end{definition}

\begin{proposition}\label{Gr1stIsOfTheFirstKind}
    The module $\mathrm{gr}^{1\mathrm{st}}(M) := M(\Gamma(M,1\mathrm{st}),U^{1\mathrm{st}},\rho^{1\mathrm{st}})$ associated to $(U^{1\mathrm{st}},\rho^{1\mathrm{st}})$ in the sense of Definition \ref{DefinitionModuleAttachedToKraftQuiver} is a twisted Gelfand-Ponomarev module of the first kind.
\end{proposition}

\begin{proof}
    First, we must prove that $FV = VF = 0$ on $\mathrm{gr}^{1\mathrm{st}}(M)$. But this is obvious by construction, since the morphisms $\rho_{E}^{1\mathrm{st}}$ are induced by the operators $F$ or $V$ on $M$ depending on whether $E$ is an $F$-arrow or a $V$-arrow respectively. Then, we must prove that all the elementary intervals of $\mathrm{gr}^{1\mathrm{st}}(M)$ are of the first kind. Given a word $w \in W$, we write $D^{1\mathrm{st}}(w)$ for the monomial induced by $w$ on $\mathrm{gr}^{1\mathrm{st}}(M)$. For $w \in W_1$, let us show by induction on $l(w)$ that 
    \begin{align}\label{ElementaryIntervalsFirstKind}
          \mathrm{Ker}D^{1\mathrm{st}}(Fw) =  U^{1\mathrm{st}}_{w} \oplus \mathrm{Dom} D^{1\mathrm{st}}(V^{\#}w).
    \end{align}
    We note that the proof is formal, that is, it does not depend on the definition of $U^{1\mathrm{st}}_w$, but only on the shape of the graph $\Gamma(M,1\mathrm{st})$ and the injectivity/surjectivity of $\rho^{1\mathrm{st}}_E$. \\

    \textbf{Step 1:} We claim that for every $w\in W_1$, $\mathrm{Dom}D^{1\mathrm{st}}(w)$ is a direct sum of $U^{1\mathrm{st}}_v$'s for $v$ running through a certain subset of $W_1$. Indeed, first notice that for every subset $I \subseteq  W_1$, we have 
    \begin{align*}
        F^{-1}\left(\bigoplus_{v \in I} U^{1\mathrm{st}}_v\right) = \sum_{v \in I} F^{-1}(U^{1\mathrm{st}}_v), & & V\left(\bigoplus_{v \in I} U^{1\mathrm{st}}_v\right) = \bigoplus_{v \in I} V(U^{1\mathrm{st}}_v).
    \end{align*}
    The right equality is obvious since any vertex of $\Gamma(M,1\mathrm{st})$ is the head of at most one $V$-arrow, and the left equality holds because 
    \begin{align*}
    \mathrm{Im}(F) = \bigoplus_{\substack{v \in W_1 \text{ s.t.}\\ vF \in W_1}} \mathrm{Im}(\rho_{vF,v}^{1\mathrm{st}}), & & \mathrm{Im}(\rho_{vF,v}^{1\mathrm{st}}) \subseteq  U^{1\mathrm{st}}_v.
    \end{align*}
    Moreover, given any $v \in W_1$, we have 
    \begin{align*}
        F^{-1}(U^{1\mathrm{st}}_v) = \begin{cases}
            \mathrm{Ker}(F) & \text{if } vF \not \in W_1;\\
            \mathrm{Ker}(F) \oplus U^{1\mathrm{st}}_{vF} & \text{if } vF \in W_1,
        \end{cases}
        & &
        V(U^{1\mathrm{st}}_v) = \begin{cases}
            0 & \text{if } vV^{\#} \not \in W_1;\\
            U^{1\mathrm{st}}_{vV^{\#}} & \text{if } vV^{\#} \in W_1.
        \end{cases}
    \end{align*}
    Here, we use the surjectivity of $\rho^{1\mathrm{st}}_E$ whenever $E = (v,vV^{\#})$ is a $V$-arrow. Besides, we have 
    \begin{equation*}
        \mathrm{Ker}(F) = U^{1\mathrm{st}}_{\emptyset} \oplus \bigoplus_{\substack{v \in W_1 \text{ s.t.} \\ vV^{\#} \in W_1}} U^{1\mathrm{st}}_{vV^{\#}}.
    \end{equation*}
    Indeed, the injectivity of $\rho^{1\mathrm{st}}_E$ whenever $E = (vF,v)$ is an $F$-arrow forces the kernel of $F$ to be reduced to the sum of the $U^{1\mathrm{st}}_v$ for $v \in W_1$ which is not the tail of any $F$-arrow. Such $v$ correspond precisely to the empty word $\emptyset$ and to all words $v \not = \emptyset$ of the first kind whose last letter is a $V^{\#}$.\\
    Now, if $w = \emptyset$ is the empty word then $\mathrm{Dom}D^{1\mathrm{st}}(w) = \mathrm{Dom}(\mathbf 1) = \mathrm{gr}^{1\mathrm{st}}(M)$ is the sum of all the $U^{1\mathrm{st}}_v$'s, so the claim holds. Assume that it holds for all $w \in W_1$ with $l(w) = k$ for some $0 \leq k < m$, and let $w \in W_1$ with $l(w) = k+1$.\\
    Assume first that $w = w'F$ for some word $w'$ which is necessarily of the first kind. Thus, we have $\mathrm{Dom}D^{1\mathrm{st}}(w) = F^{-1}(\mathrm{Dom}D^{1\mathrm{st}}(w'))$. By induction, $\mathrm{Dom}D^{1\mathrm{st}}(w')$ is a direct sum of certain $U^{1\mathrm{st}}_v$'s. By our discussion above, it follows that $\mathrm{Dom}D^{1\mathrm{st}}(w)$ itself can be written as such a direct sum.\\
    \sloppy Assume now that $w = w'V^{\#}$ for some $w' \in W_1$. We have $\mathrm{Dom}D^{1\mathrm{st}}(w) = (V^{\#})^{-1}(\mathrm{Dom}D^{1\mathrm{st}}(w')) = V(\mathrm{Dom}D^{1\mathrm{st}}(w'))$. Again, by induction and the discussion above, the claim follows.\\

    \textbf{Step 2:} We prove \eqref{ElementaryIntervalsFirstKind} by induction on $l(w)$. First, observe that
    \begin{equation*}
        \mathrm{Dom}(V^{\#}) = \mathrm{Im}(V) = \bigoplus_{\substack{v \in W_1 \text{ s.t.} \\ vV^{\#} \in W_1}} U^{1\mathrm{st}}_{vV^{\#}}.
    \end{equation*}
     In particular, we see that 
    \begin{equation*}
        \mathrm{Ker}(F) =  U^{1\mathrm{st}}_{\emptyset} \oplus \mathrm{Dom}(V^{\#}),
    \end{equation*}
    which is precisely \eqref{ElementaryIntervalsFirstKind} for $w$ equal to the empty word. Now, let us assume that \eqref{ElementaryIntervalsFirstKind} is known for all $w \in W_1$ of length $k$ for some $0 \leq k < m$, and let $w \in W_1$ have length $k+1$. Assume first that $w = w'F$ for some word $w' \in W_1$. By induction, we have
    \begin{equation*}
        \mathrm{Ker}D^{1\mathrm{st}}(Fw) = F^{-1}(\mathrm{Ker}D^{1\mathrm{st}}(Fw')) = F^{-1}\left( U^{1\mathrm{st}}_{w'} \oplus \mathrm{Dom} D^{1\mathrm{st}}(V^{\#}w')\right).
    \end{equation*}
    By step 1, $F^{-1}$ preserves the sum on the right-hand side so that we have 
    \begin{equation*}
        \mathrm{Ker}D^{1\mathrm{st}}(Fw) = F^{-1}(U^{1\mathrm{st}}_{w'}) + F^{-1}(\mathrm{Dom} D^{1\mathrm{st}}(V^{\#}w')) = U_{w}^{1\mathrm{st}} + \mathrm{Dom}D^{1\mathrm{st}}(V^{\#}w).
    \end{equation*}
    \sloppy Here, we used the fact that $\mathrm{Dom}D^{1\mathrm{st}}(V^{\#}w)$ contains $\mathrm{Ker}(F)$. It remains to show that the sum is direct. So let $x \in U_{w}^{1\mathrm{st}} \cap \mathrm{Dom}D^{1\mathrm{st}}(V^{\#}w)$. Then $\rho_{w,w'}^{1\mathrm{st}}(x) \in U_{w'}^{1\mathrm{st}} \cap \mathrm{Dom} D^{1\mathrm{st}}(V^{\#}w') = \{0\}$ by induction. Since $\rho_{w,w'}^{1\mathrm{st}}$ is injective, we have $x = 0$ as desired. \\
    Assume now that $w = w'V^{\#}$ for some word $w' \in W_1$. By induction we have
    \begin{equation*}
        \mathrm{Ker}D^{1\mathrm{st}}(Fw) = V(\mathrm{Ker}D^{1\mathrm{st}}(Fw')) = V\left( U^{1\mathrm{st}}_{w'} \oplus \mathrm{Dom} D^{1\mathrm{st}}(V^{\#}w')\right).
    \end{equation*}
    By Step 1, $\mathrm{Dom} D^{1\mathrm{st}}(V^{\#}w')$ is a direct sum of some $U^{1\mathrm{st}}_v$'s. By induction, this sum does not involve $U^{1\mathrm{st}}_{w'}$. Thus, we have 
    \begin{equation*}
        \mathrm{Ker}D^{1\mathrm{st}}(Fw) = V(U^{1\mathrm{st}}_{w'}) \oplus V(\mathrm{Dom} D^{1\mathrm{st}}(V^{\#}w')) = U_{w}^{1\mathrm{st}} + \mathrm{Dom}D^{1\mathrm{st}}(V^{\#}w),
    \end{equation*}
    which concludes the proof of Step 2. \\

    Since $U^{1\mathrm{st}}_{w} \not = \{0\}$ for all $w \in W_1$, we deduce that all the monomials of the form $D^{1\mathrm{st}}(w)$ for $w \in W_1$ are monomials of the first kind in $\mathrm{gr}^{1\mathrm{st}}(M)$. Besides, if $w \in W_1$ then the corresponding elementary interval $(\beta_i^{1\mathrm{st}}, \beta_{i+1}^{1\mathrm{st}})$ of the first kind in $\mathrm{gr}^{1\mathrm{st}}(M)$ satisfies 
    \begin{equation*}
        \beta_{i+1}^{1\mathrm{st}}/\beta_i^{1\mathrm{st}} = U^{1\mathrm{st}}_{w}.
    \end{equation*}
    Since $\mathrm{gr}^{1\mathrm{st}}(M) = \bigoplus_{w \in W_1} U^{1\mathrm{st}}_{w}$, by dimensions we have already exhausted all the elementary intervals of $\mathrm{gr}^{1\mathrm{st}}(M)$. This concludes the proof.
\end{proof}

\end{document}
